\documentclass[11pt]{article}
\usepackage{amsmath,amssymb,amsthm}
\usepackage[margin=1.2in]{geometry}
\newtheorem{theorem}{Theorem}
\newtheorem{lemma}[theorem]{Lemma}
\newtheorem{proposition}[theorem]{Proposition}
\newtheorem{corollary}[theorem]{Corollary}
\newtheorem{remark}[theorem]{Remark}
\title{Erd\H{o}s \#400: the ledger, a short proof of the sharp upper bound,\\
and the constant}
\author{Samuel Lavery}
\date{Draft --- \today}
\begin{document}
\maketitle

\begin{abstract}
For $k\ge2$ let $g_k(n)=\max\{(a_1+\cdots+a_k)-n:\ a_1!\cdots a_k!\mid n!\}$.  We record an
exact reformulation of the divisibility as a per-rail digit-sum ledger, and from it give a
five-line unconditional proof of
$g_k(n)\le(k-1)\log_2 n+\log_2\log n+O_k(1)$, matching the upper bound of Li
\cite{Li400} --- the bound comes from the prime $2$ alone, because the ledger charges
$p-1$ per carry and $2$ is the only rail whose cell is trivial.  We then give an exact
polylogarithmic algorithm for the distribution of the carry count on each rail, validated
against exhaustive computation, and use it to measure the constant of Erd\H{o}s and Graham
out to $n=10^{31}$.  The measurement says the upper bound is \emph{sharp}:
\[
  \sum_{n\le x}g_k(n)\sim c_k\,x\log x\quad\text{with}\quad
  c_k=\frac{k-1}{\log 2}=1.44270\ldots(k-1),
\]
against Li's proven bracket $c_k\in(k-1)[\,3/\log12,\ 1/\log2\,]=(k-1)[1.2073,1.4427]$.
The constant is \textbf{measured, not proved}; the upper bound and the ledger are proved.
\end{abstract}

\section{The ledger}

Write $s_p(m)$ for the sum of base-$p$ digits of $m$.  Let $c_j$ denote the carry out of
position $j$ in the base-$p$ addition $a_1+\cdots+a_k$, and set
\[
  \kappa_p(a_1,\dots,a_k)\;:=\;\sum_{j\ge0}c_j ,
\]
the \emph{total carry weight}.  Throughout $A=a_1+\cdots+a_k$ and $G=A-n$.

\begin{remark}[weight, not count --- a correction]\label{rem:weight}
For $k=2$ every $c_j\in\{0,1\}$ and $\kappa_p$ is the \emph{number} of carrying positions,
which is Kummer's theorem in its usual form.  For $k\ge3$ a carry may exceed $1$ (up to
$k-1$, by Lemma~\ref{lem:cap}), and the identity below holds for the total weight
$\sum_jc_j$, \emph{not} for the count of nonzero carries.  An earlier draft of this note
said ``number of carries'' throughout; that is correct only at $k=2$.  Caught in
formalisation.
\end{remark}

\begin{theorem}[ledger identity]\label{thm:ledger}
For non-negative integers $a_1,\dots,a_k$ and $n$,
\[
  a_1!\cdots a_k!\ \Big|\ n!
  \qquad\Longleftrightarrow\qquad
  A-n\ \le\ \sum_{i=1}^k s_p(a_i)-s_p(n)\quad\text{for every prime }p.
\]
Equivalently, since $\sum_i s_p(a_i)-s_p(A)=(p-1)\kappa_p$,
\[
  g_k(n)=\max_{a}\ \min_{p}\ \Bigl[(p-1)\,\kappa_p(a)\ +\ s_p(n+G)-s_p(n)\Bigr].
\]
\end{theorem}

\begin{proof}
By Legendre, $v_p(m!)=\frac{m-s_p(m)}{p-1}$, so
$v_p(n!)-\sum_iv_p(a_i!)=\frac{(n-s_p(n))-\sum_i(a_i-s_p(a_i))}{p-1}
=\frac{\sum_is_p(a_i)-s_p(n)-(A-n)}{p-1}$.
Divisibility is exactly non-negativity of this for all $p$.  The second form is Kummer's
theorem $\sum_is_p(a_i)-s_p(A)=(p-1)\kappa_p$.
\end{proof}

\emph{Verified}: $0$ mismatches against brute-force factorial divisibility for all
$n<40$, $a_1\le n+8$, $G\le4$, with $p$ ranging to $2n+8$.  (Truncating the rail set is
unsafe: with $p<60$ only, $g_2$ is already wrong for $10$ values of $n\le400$.)

Two features of the ledger drive everything.  First, \emph{a carry on rail $p$ pays
$p-1$}, so the smallest rails bind: in computation $p=2$ binds $73\%$ of the time, $p=3$
a further $20\%$.  Second, every rail with $p>\max_i a_i$ contributes exactly $G$, so
large rails are \emph{always exactly tight} and the binding structure lives entirely on
small rails.

\section{The upper bound}

\begin{lemma}[universal carry cap]\label{lem:cap}
In the base-$p$ addition of $k$ non-negative integers, the carry out of every position is
at most $k-1$.  Consequently
$\kappa_p(a_1,\dots,a_k)\le(k-1)\bigl(\lfloor\log_pA\rfloor+1\bigr)$.
\end{lemma}

\begin{proof}
Let $c_j$ be the carry out of position $j$, $c_{-1}=0$.  Then
$c_j=\bigl\lfloor(\sum_i a_{i,j}+c_{j-1})/p\bigr\rfloor
\le\bigl\lfloor(k(p-1)+c_{j-1})/p\bigr\rfloor$.
If $c_{j-1}\le k-1$ then the right side is at most
$\lfloor(kp-k+k-1)/p\rfloor=\lfloor(kp-1)/p\rfloor=k-1$; induction gives $c_j\le k-1$ for
all $j$.  Positions beyond $\lfloor\log_pA\rfloor$ contribute no carry.
\end{proof}

\emph{Verified}: $40\,000$ random $(k,p,a_1,\dots,a_k)$ with $k\le6$, $p\le11$,
$a_i<10^8$: no violation of either statement.

\begin{theorem}[sharp upper bound]\label{thm:upper}
For every $k\ge2$ and all large $n$,
\[
  g_k(n)\ \le\ (k-1)\log_2 n+\log_2\log_2 n+O_k(1),
  \qquad\text{hence}\qquad
  c_k\le\frac{k-1}{\log 2}.
\]
\end{theorem}

\begin{proof}
Take $p=2$ in Theorem~\ref{thm:ledger}.  Since $p-1=1$,
\[
  G\ \le\ \kappa_2+s_2(A)-s_2(n).
\]
By Lemma~\ref{lem:cap}, $\kappa_2\le(k-1)(\lfloor\log_2A\rfloor+1)$.  By subadditivity of
the binary digit sum, $s_2(A)=s_2(n+G)\le s_2(n)+s_2(G)$, so
$s_2(A)-s_2(n)\le s_2(G)\le\log_2(G+1)$.  Therefore
\begin{equation}\label{eq:discrete}
  G\ \le\ (k-1)\log_2A+\log_2(G+1)+O(k).
\end{equation}
It remains to remove the $G$ on the right, and \eqref{eq:discrete} does this by itself.
Write $C=(k-1)\log_2A+O(k)$, so that $G\le C+\log_2(G+1)$.  Since $G\le A$ we get the crude
bound $G\le C+\log_2(A+1)$, hence $G=O_k(\log A)$; feeding that back into
\eqref{eq:discrete} gives $\log_2(G+1)\le\log_2\log_2A+O_k(1)$ and therefore
\[
  G\ \le\ (k-1)\log_2A+\log_2\log_2A+O_k(1).
\]
Finally $\log_2A=\log_2n+O(1/n)$ because $A=n+G$ with $G=O_k(\log n)$.
\end{proof}

\begin{remark}[no external input is needed]\label{rem:selfcontained}
An earlier draft absorbed the $\log_2(G+1)$ term by quoting Erd\H{o}s and Graham's bound
$g_k(n)\ll_k\log n$.  That citation is unnecessary: as above, \eqref{eq:discrete} bootstraps
against itself in one step, since $G\le A$ already gives $G=O_k(\log A)$.  So
Theorem~\ref{thm:upper} is self-contained, depending only on Legendre, Kummer and
Lemma~\ref{lem:cap}.  Noticed in formalisation, and the formal proof takes this route.
\end{remark}

\begin{remark}
The bound is forced by a single rail, and by the \emph{only} rail whose cell $p-1$ equals
$1$: for general $p$ the same argument gives
$G\le(p-1)(k-1)\log_pA+O(\log\log n)$, and $\min_p\frac{(p-1)}{\log p}$ is attained at
$p=2$ with value $1/\log2$.  Deleting rail $2$ from the constraint set raises the measured
constant by a factor $1.35$; restricting to $p\equiv1\ (\mathrm{mod}\ 6)$ raises it by
$2.95$.
\end{remark}

\begin{remark}[the cap is attained]
The carry recursion $c=\lfloor(k+c)/2\rfloor$ has fixed point $c=k-1$, realised when every
binary digit of every $a_i$ equals $1$, i.e.\ $a_i=2^{m}-1$.  So Lemma~\ref{lem:cap} is
sharp on rail $2$, and the upper bound of Theorem~\ref{thm:upper} is attainable
\emph{as a constraint}.  Whether such splits are available for a given $n$, subject to the
remaining rails, is the lower-bound problem; the best proved lower bound is Li's
$\bigl(\frac{3(k-1)}{\log12}-\varepsilon\bigr)\log n$ for density-one $n$ \cite{Li400}.
\end{remark}

\section{Monotonicity in the number of parts}

The upper bound of \S2 is complemented by an exact monotonicity, proved by appending a
trivial part.

\begin{theorem}[append a one]\label{thm:append}
For all $n\ge1$ and all $k\ge3$,
\[
  g_k(n)\ \ge\ g_{k-1}(n)+1 ,
\]
and consequently $g_k(n)\ge g_2(n)+k-2\ge k-1$ for every $k\ge2$.
\end{theorem}

\begin{proof}
Let $b_1,\dots,b_{k-1}$ attain $g_{k-1}(n)$, so $b_1!\cdots b_{k-1}!\mid n!$ and
$\sum_{i<k}b_i=n+g_{k-1}(n)$.  Put $a_i=b_i$ for $i\le k-1$ and $a_k=1$.  Then
$\prod_{i\le k}a_i!=\prod_{i<k}b_i!$ divides $n!$, so the $k$-tuple is admissible, while
$\sum_{i\le k}a_i-n=g_{k-1}(n)+1$.  As $g_k(n)$ is a supremum over admissible $k$-tuples,
the inequality follows.  Iterating down to $k=2$ gives $g_k(n)\ge g_2(n)+k-2$; and
$b_1=n$, $b_2=1$ is admissible with $b_1+b_2-n=1$, so $g_2(n)\ge1$.
\end{proof}

\begin{remark}[why the step is exactly tight]\label{rem:tight}
Read on the ledger of Theorem~\ref{thm:ledger}, appending $a_k=1$ raises $s_p(a_k)$ by
exactly $1$ on \emph{every} rail simultaneously, while $G$ also rises by exactly $1$.  So
the constraint $G\le\min_p\bigl[\sum_is_p(a_i)-s_p(n)\bigr]$ is preserved with equality in
the increment, rail by rail: no rail is left slack and none is violated.  This is why $1$
is the correct part to append and why the step cannot be improved by another choice of
$a_k$ --- any $a_k\ge2$ raises $G$ by $a_k$ but raises $s_p(a_k)$ by only $a_k$ for
$p>a_k$ and by strictly less for $p\le a_k$, so some small rail loses.
\end{remark}

\begin{corollary}\label{cor:ck}
$c_k\ge c_2$ for every $k\ge2$, so with Theorem~\ref{thm:upper},
$c_2\le c_k\le\frac{k-1}{\log2}$.  In particular the $k$-dependence of Erd\H{o}s and
Graham's $g_k(n)\ll_k\log n$ cannot be made uniform in $k$.
\end{corollary}

\begin{remark}[scope of the append step]\label{rem:scope}
The move applies exactly when the quantity is a supremum over tuples whose admissibility is
a divisibility between products of factorials, since then $1!=1$ is a free unit.  It
therefore transfers to Erd\H{o}s \#391 (largest least part in a factorisation $n!=a_1\cdots
a_n$), giving monotonicity in the number of parts by the same argument.  It applies only
one-sidedly to \#374 ($a_1<\cdots<a_k=m$ with $a_1!\cdots a_k!$ a square): appending $1$
preserves squareness because $w(1)=0$ in the $\mathbb F_2$ valuation ledger, but strict
increase permits it only at the bottom and only when $a_1>1$.  It does not apply to
\#683, \#685, \#700, \#727, \#1093, \#1094, \#1095, where the quantity is a least $n$, an
extreme prime factor, a divisor count or a gcd minimum: those load every rail onto a single
carrier point, so there is no tuple coordinate to extend.
\end{remark}

\section{An exact algorithm, and the constant}

\begin{proposition}[carry automaton]\label{prop:dp}
For fixed $A$ and $p$, the exact distribution of $\kappa_p$ over all splits
$a_1+a_2=A$ is computed in $O(\log^2_pA)$ arithmetic operations by dynamic programming on
$(\text{position},\ \text{carry},\ \#\text{carries})$: with $A_j$ the $j$-th digit of $A$,
a transition from carry-in $c$ to carry-out $c'$ at position $j$ is weighted by the number
of pairs $(x,y)\in[0,p-1]^2$ with $x+y=T:=A_j+pc'-c$, namely $T+1$ if $T\le p-1$ and
$2p-1-T$ otherwise.
\end{proposition}

\emph{Verified} against brute force for all $A<180$ and $p\in\{2,3,5,7\}$.

\begin{remark}[attribution: this is Holte's carries matrix]\label{rem:holte}
The unconditioned version of this automaton --- the Markov chain on carries when $k$ uniform
base-$p$ digits are added --- is not new.  It is the ``amazing matrix'' of Holte, \emph{Carries,
combinatorics, and an amazing matrix}, Amer.\ Math.\ Monthly \textbf{104} (1997), 138--149,
revisited by Diaconis and Fulman, \emph{Carries, shuffling, and an amazing matrix}, Amer.\ Math.\
Monthly \textbf{116} (2009), 788--803.  Two of its properties are worth importing here, and both
were verified directly for $k\in\{2,3,4\}$ and $p\in\{2,3,5\}$ to machine zero:
\begin{enumerate}
\item[(i)] the transition matrix has eigenvalues exactly
      $1,\ p^{-1},\ p^{-2},\ \dots,\ p^{-(k-1)}$;
\item[(ii)] its stationary distribution is the Eulerian distribution
      $\pi(j)=A(k,j)/k!$ on $j=0,\dots,k-1$ --- and is \emph{independent of $p$}
      ($\tfrac12,\tfrac12$ at $k=2$; $\tfrac16,\tfrac23,\tfrac16$ at $k=3$;
      $\tfrac1{24},\tfrac{11}{24},\tfrac{11}{24},\tfrac1{24}$ at $k=4$).
\end{enumerate}
So the \emph{typical} carry per position is Eulerian with mean $(k-1)/2$, in closed form and
with no dynamic programming.  What Proposition~\ref{prop:dp} adds is the conditioning on a
prescribed sum $A$, which the classical chain does not carry; the digit-pair weight
$T+1$ or $2p-1-T$ is the conditioned transition.  The distinction matters for
$g_k$: the constant $c_k=(k-1)/\log2$ corresponds to \emph{saturating} every base-$2$
position at the maximal carry $k-1$, not to the Eulerian mean $(k-1)/2$.  The DC of this
problem is therefore known exactly and is not what $g_k$ measures --- $g_k$ is a maximum over
tuples, which is why the extremal reading is the right one here.
\end{remark}

Writing $R_p=\lceil(G-(s_p(A)-s_p(n)))/(p-1)\rceil$ for the carry demand on rail $p$ and
$\mathrm{surv}_p(G)$ for the DP probability that $\kappa_p\ge R_p$, the first-moment
estimator is
\[
  \widehat g_k(n)=\max\Bigl\{G:\ (A+1)\prod_p\mathrm{surv}_p(G)\ \ge\ 1\Bigr\}.
\]
Two facts license it.  The cross-rail correlation factor between the true joint survival
and the product is $\Gamma=1.03\pm0.34$ over $n\le1200$ --- the rails are effectively
independent.  And $\widehat g_2$ is \emph{unbiased} against exhaustive $g_2$ on
$n\in[300,1600]$: mean bias $-0.05$, ratio $0.9943$, agreement within $\pm1$ pointwise.

\emph{Measured.}  With $\widehat g_2$ evaluated on random $n$ at each scale
($\log n$ from $10$ to $70$):
\[
\begin{array}{r|cccccc}
 n & 10^{12} & 10^{15} & 10^{18} & 10^{22} & 10^{26} & 10^{30}\\\hline
 \widehat g_2/\log n & 1.4022 & 1.4070 & 1.4136 & 1.4325 & 1.4352 & 1.4229\\
 \text{gap to }1/\log2 & -0.041 & -0.036 & -0.029 & -0.010 & -0.008 & -0.020
\end{array}
\]
The gap to $1/\log2=1.44270$ shrinks; the gap to $3/\log12=1.20729$ stays near $+0.21$
throughout.  Hence the prediction
\[
  \boxed{\ c_k=\frac{k-1}{\log 2}\ }
\]
i.e.\ Li's upper bound is sharp and his density-one lower bound is not the truth.

\begin{remark}[why earlier attempts missed it]
Two errors, both instructive.  Estimating $c_2$ by averaging per-$n$ maxima drifts and does
not converge; averaging the occupancy \emph{first} and taking the quantile of the mean is
stable.  And over $n\le10^4$ the functions $\log_2n$ and $\log_2\log n$ are nearly
collinear, so a fit on that range cannot separate $1.2073$ from $1.4427$; the carry
automaton removes the enumeration cost and lengthens the lever from
$\log n\in[5.6,8.8]$ to $[10,70]$, at which point the separation is unambiguous.
\end{remark}

\begin{remark}[register]
\textbf{Proved}: Theorem~\ref{thm:ledger}, Lemma~\ref{lem:cap},
Theorem~\ref{thm:upper}, Theorem~\ref{thm:append}, Corollary~\ref{cor:ck},
Proposition~\ref{prop:dp}.  Theorem~\ref{thm:upper} recovers Li's
upper bound by a different and shorter route; no priority is claimed for it.
Theorem~\ref{thm:append} is elementary and gives $c_2\le c_k\le(k-1)/\log2$.
\textbf{Measured}: $c_k=(k-1)/\log2$, with the three validations above.
\textbf{Not proved}: the lower bound, hence \#400 itself.
\textbf{Falsifier, pre-registered}: $\widehat g_2/\log n$ settling below $1.40$ at
$n\ge10^{40}$, or $\widehat g_2$ developing bias against exhaustive $g_2$.
\end{remark}

\begin{thebibliography}{9}
\bibitem{Li400} E.\ Li, \emph{Prime-power rarefaction and a density-one lower bound for
Erd\H{o}s problem 400}, arXiv:2606.23661v2 (2026).
\bibitem{ErGr80} P.\ Erd\H{o}s and R.\ L.\ Graham, \emph{Old and new problems and results
in combinatorial number theory}, Monogr.\ Enseign.\ Math.\ 28 (1980), p.~77.
\end{thebibliography}
\end{document}
